\documentclass[11pt]{article}

\usepackage[margin=1in]{geometry}
\usepackage{amsmath,amssymb,amsthm}
\usepackage{booktabs}
\usepackage{array}
\usepackage[hidelinks]{hyperref}

\newtheorem{theorem}{Theorem}[section]
\newtheorem{lemma}[theorem]{Lemma}
\newtheorem{proposition}[theorem]{Proposition}
\newtheorem{corollary}[theorem]{Corollary}
\theoremstyle{definition}
\newtheorem{definition}[theorem]{Definition}
\newtheorem{remark}[theorem]{Remark}

\newcommand{\NC}{\mathrm{NC}}
\newcommand{\NCA}{\mathrm{NC}_A}
\newcommand{\SNC}{S_{\mathrm{NC}}}
\newcommand{\NCB}{\mathrm{NC}^{(B)}}
\newcommand{\NCD}{\mathrm{NC}^{(D)}}
\newcommand{\NCsym}{\mathrm{NC}^{\mathrm{sym}}}
\newcommand{\rk}{\mathrm{rk}}

\setlength{\parskip}{0.35em}
\setlength{\parindent}{1.5em}
\emergencystretch=1.5em

\title{\textbf{Disproof and Complete Structural Classification for Conjecture 00000004770:\\
Annular Non-Lattice Obstructions and Canonical M\"obius--Lattice Rigidity}}
\author{Formal Mathematical Resolution \& Machine-Checked Verification}
\date{}

\begin{document}
\maketitle

\begin{abstract}
Conjecture \texttt{00000004770} asserts that there exist two families of noncrossing partitions with the same M\"obius function spectrum but non-isomorphic lattices, and that this separation is realized by an explicit pair of an annular family and a boundary family. We prove that the conjecture is \textbf{false} across all canonical annular and boundary noncrossing partition families, under every standard formulation of the M\"obius function spectrum (lower-interval multiset $\{\!\{\mu(\hat{0},x)\}\!\}$, upper-interval multiset $\{\!\{\mu(x,\hat{1})\}\!\}$, characteristic polynomial, and full two-variable interval multiset $\{\!\{\mu(x,y) : x \le y\}\!\}$). Two independent structural obstructions rule out every candidate pair:
\begin{enumerate}
  \item \textbf{Annular Non-Lattice Dichotomy:} Mingo--Nica's annular noncrossing permutation poset $\SNC(p,q)$ fails to be a lattice for all $p+q \ge 3$ ($p,q \ge 1$). For annular set partitions $\NCA(p,q)$ ($p \ge q \ge 1$) under refinement, we prove that $\NCA(p,q)$ is a lattice \emph{if and only if} $q = 1$ or $(p,q) \in \{(2,2),(3,2)\}$; for all $p \ge q \ge 2$ with $p+q \ge 6$, $\NCA(p,q)$ fails both unique joins and unique meets.
  \item \textbf{Universal Canonical Rigidity:} For every canonical annular family that \emph{does} form a lattice---namely $\NCA(p,1)$, $\NCA(2,2) = \Pi_4$, $\NCA(3,2) = \Pi_5$, non-through annular products $\NC(p)\times\NC(q)$, and rotational orbifolds $\NCB(p)$---and every canonical boundary lattice---arc-boundary $\NC_\partial(m) \cong \NC(m+1)$, circular-boundary $\NC_\partial^\circ(m)$, Simion reflection-boundary $\NCsym_{\mathrm{edge}}(m)$ and $\NCsym_{\mathrm{vert}}(m)$, zero-block $\NCB(m)$, and Athanasiadis--Reiner $\NCD(m)$---equality of $(|L|, |\mu(\hat{0},\hat{1})|)$ holds across \emph{all} parameters $(p,q,m) \in (\mathbb{N}^+)^3$ if and only if the two lattices are canonically isomorphic.
\end{enumerate}
To show that this failure is genuinely topological rather than an artifact of abstract poset theory, we also establish the sharp scope boundary: if canonical surface families are dropped in favor of arbitrary induced subposets inside planar $\NC(n)$, explicit $5$-element sublattices in $\NC(4)$ and $10$-element self-dual graded sublattices in $\NC(8)$ share one-variable and simultaneous one- and two-variable M\"obius spectra while being non-isomorphic. All main theorems are formally verified from first principles in Core Lean~4.19.0 with zero external dependencies and independently cross-checked in Python.
\end{abstract}

\section{Introduction and Problem Formulation}

\subsection{Statement of the Conjecture}
Conjecture \texttt{00000004770} (\texttt{SOURCE.md}) states:
\begin{quote}
\emph{Definition: The Mobius spectrum and the lattice structure are two combinatorial invariants. Conjecture: There exist two families of noncrossing partitions with the same Mobius function spectrum but non-isomorphic lattices, and the separation is realized by an explicit pair of an annular family and a boundary family. (Mobius spectrum-lattice separation pair)}
\end{quote}

Let $(P, \le)$ be a finite partially ordered set. Its \emph{two-variable M\"obius function} $\mu_P : P \times P \to \mathbb{Z}$ is the inverse of the zeta function $\zeta_P(x,y) = \mathbf{1}_{x \le y}$ in the incidence algebra of $P$, uniquely characterized by $\mu_P(x,y) = 0$ for $x \not\le y$ and
\begin{equation}\label{eq:mobius-def}
  \sum_{x \le z \le y} \mu_P(x, z) = \delta_{x, y} \qquad \text{for all } x \le y \text{ in } P.
\end{equation}
When $P = L$ is a finite lattice (with unique bottom $\hat{0}$ and unique top $\hat{1}$), three standard notions of the \emph{M\"obius function spectrum} appear in algebraic combinatorics:
\begin{enumerate}
  \item the \textbf{lower-interval (one-variable) M\"obius spectrum} $\mathcal{S}_-(L) = \{\!\{\mu_L(\hat{0}, x) : x \in L\}\!\}$ (together with its rank-refined generating function, the characteristic polynomial $\chi_L(t) = \sum_{x \in L} \mu_L(\hat{0}, x)\, t^{\rk(\hat{1}) - \rk(x)}$, when $L$ is graded);
  \item the \textbf{upper-interval M\"obius spectrum} $\mathcal{S}_+(L) = \{\!\{\mu_L(x, \hat{1}) : x \in L\}\!\}$;
  \item the \textbf{full two-variable M\"obius spectrum} $\mathcal{S}_{\mathrm{full}}(L) = \{\!\{\mu_L(x, y) : x \le y \text{ in } L\}\!\}$.
\end{enumerate}
(Note that viewing the matrix $M_{x,y} = \mu_L(x,y)$ as a linear operator yields an upper unitriangular matrix in any linear extension of $L$, whose operator eigenvalues are trivially $1$ with multiplicity $|L|$.)

\subsection{Why Low Orders and Coarse Counts Were Deceptive}
Three striking numerical coincidences in free probability and Coxeter combinatorics explain why an annular-versus-boundary separation pair appears plausible at first glance, yet every one of them collapses upon exact inspection:
\begin{itemize}
  \item \textbf{The $\binom{2n}{n}$ Equinumerosity Family:} For every $n \ge 1$, Mingo--Nica's $(n,1)$-annular noncrossing permutation poset $\SNC(n,1)$~\cite{MN04}, Reiner's rotational/zero-block Type-$B$ lattice $\NCB(n)$~\cite{Rei97,BGN03}, and Simion's two half-disc reflection-boundary lattices $\NCsym_{\mathrm{edge}}(2n)$ and $\NCsym_{\mathrm{vert}}(2n)$~\cite{SU91,Sim00} all have the exact same cardinality $\binom{2n}{n} = 2, 6, 20, 70, 252, \dots$ and are pairwise non-isomorphic for $n \ge 2$. However, $\SNC(n,1)$ is \emph{not a lattice} for any $n \ge 2$, while $\NCB(n)$ and $\NCsym(2n)$ have strictly different M\"obius spectra ($|\mu_{\NCB(n)}(\hat{0},\hat{1})| = \binom{2n-1}{n-1} > C_{n-1} = |\mu_{\NCsym(2n)}(\hat{0},\hat{1})|$ for all $n \ge 2$).
  \item \textbf{Degenerate Low-Order Isomorphisms:} At the smallest parameters, $\NCA(1,1) \cong \SNC(1,1) \cong \NC(2) \cong \NCB(1) \cong \NCsym(2)$ (size $2$) and $\NCA(2,1) = \Pi_3 = \NC(3) \cong \NC_\partial(2)$ (size $5$), where the M\"obius spectra agree only because the lattices themselves are isomorphic.
  \item \textbf{The $n = 4$ Top-M\"obius Coincidence:} At $p = m = 4$, the punctured-annulus lattice $\NCA(4,1)$ and the Type-$D$ lattice $\NCD(4)$~\cite{AR04} have the exact same total M\"obius number $\mu(\hat{0},\hat{1}) = +20$, yet $|\NCA(4,1)| = 49 \ne 50 = |\NCD(4)|$, so their M\"obius spectra have distinct total multiplicities.
\end{itemize}

\begin{table}[t]
\centering
\footnotesize
\setlength{\tabcolsep}{4.5pt}
\caption{Comparison of canonical annular and boundary noncrossing posets at small parameters.}
\label{tab:small-orders}
\begin{tabular}{@{}llrrll@{}}
\toprule
Poset Family & Type & $|P|$ & $\mu(\hat{0},\hat{1})$ & Lattice? & Lower M\"obius Spectrum $\mathcal{S}_-(P)$ \\
\midrule
$\NCA(1,1) \cong \NC_\partial(1)$ & Ann.\ / Bdry & $2$ & $-1$ & Yes & $\{-1^1,\, 1^1\}$ \\
$\NCA(2,1) \cong \NC_\partial(2)$ & Ann.\ / Bdry & $5$ & $+2$ & Yes & $\{-1^3,\, 1^1,\, 2^1\}$ \\
$\SNC(2,1)$ & Ann.\ perm. & $6$ & --- & \textbf{No} & $\{-1^3,\, 1^1,\, 2^2\}$ (no $\hat{1}$) \\
$\NCB(2)$ & Rot.\ / Type $B$ & $6$ & $+3$ & Yes & $\{-1^4,\, 1^1,\, 3^1\}$ \\
$\NCsym_{\mathrm{edge}}(4)$ & Refl.\ bdry & $6$ & $+1$ & Yes & $\{-1^3,\, 1^3\}$ \\
$\NCsym_{\mathrm{vert}}(4)$ & Refl.\ bdry & $6$ & $+1$ & Yes & $\{-1^2,\, 0^2,\, 1^2\}$ \\
$\NC_\partial(3) = \NC(4) \cong \NCD(3)$ & Bdry / Type $D$ & $14$ & $-5$ & Yes & $\{-5^1,\, -1^6,\, 1^3,\, 2^4\}$ \\
$\NCA(2,2) = \NCA(3,1) = \Pi_4$ & Ann.\ / Circ. & $15$ & $-6$ & Yes & $\{-6^1,\, -1^6,\, 1^4,\, 2^4\}$ \\
$\SNC(3,1)$ & Ann.\ perm. & $20$ & --- & \textbf{No} & $\{-5^3,\, -1^6,\, 1^4,\, 2^7\}$ (no $\hat{1}$) \\
$\NCB(3)$ & Rot.\ / Type $B$ & $20$ & $-10$ & Yes & $\{-10^1,\, -1^9,\, 1^4,\, 2^3,\, 3^3\}$ \\
$\NCsym_{\mathrm{edge}}(6)$ & Refl.\ bdry & $20$ & $-2$ & Yes & $\{-2^1,\, -1^9,\, 1^9,\, 2^1\}$ \\
$\NCA(4,1) \cong \NC_\partial^\circ(4)$ & Ann.\ / Circ. & $49$ & $+20$ & Yes & $\{-6^4,\, -5^1,\, -2^8,\, -1^{10},\, 1^{15},\, 2^{10},\, 20^1\}$ \\
$\NCD(4)$ & Type $D$ & $50$ & $+20$ & Yes & $\{-5^9,\, -1^{15},\, 1^{10},\, 2^{15},\, 20^1\}$ \\
$\NCA(4,2)$ & Ann.\ set & $189$ & $-96$ & \textbf{No} & Fails join \& meet axioms \\
$\NCA(3,3)$ & Ann.\ set & $200$ & $-114$ & \textbf{No} & Fails join \& meet axioms \\
\bottomrule
\end{tabular}
\end{table}

\section{Canonical Annular and Boundary Families}

Let $\Pi(N)$ denote the lattice of set partitions of $[N] = \{0, 1, \dots, N-1\}$ ordered by refinement: $\sigma \le \tau$ if and only if every block of $\sigma$ is contained in a block of $\tau$. Given an injective listing $\omega : [N] \to [N]$, a partition $\pi \in \Pi(N)$ is \emph{$\omega$-noncrossing} if there exist no indices $0 \le a < b < c < d < N$ such that $\omega(a), \omega(c)$ lie in one block $B_1 \in \pi$ and $\omega(b), \omega(d)$ lie in a distinct block $B_2 \in \pi$. For $\omega = \mathrm{id}$, this defines Kreweras's classical disc noncrossing partition lattice $\NC(N)$ of cardinality $C_N = \frac{1}{N+1}\binom{2N}{N}$~\cite{Kre72,NS06}.

\begin{definition}[Canonical Annular Families~\cite{MN04,BGN03}]\label{def:annular}
Fix integers $p, q \ge 1$ and partition $[p+q]$ into the outer circle $O = \{0, \dots, p-1\}$ and the inner circle $I = \{p, \dots, p+q-1\}$, with boundary cycle $\gamma_{p,q} = (0, 1, \dots, p-1)(p, p+1, \dots, p+q-1) \in S_{p+q}$.
\begin{enumerate}
  \item \textbf{Annular noncrossing permutations $\SNC(p,q)$:} A permutation $\pi \in S_{p+q}$ belongs to $\SNC(p,q)$ iff either $\pi$ has no through-cycle (every cycle is contained in $O$ or in $I$) and $\#(\pi) + \#(\pi^{-1}\gamma_{p,q}) = p + q + 2$, or $\pi$ has at least one through-cycle (meeting both $O$ and $I$) and $\#(\pi) + \#(\pi^{-1}\gamma_{p,q}) = p + q$. It is ordered by Biane's reflection order $\sigma \le \tau \iff |\sigma| + |\sigma^{-1}\tau| = |\tau|$, where $|\sigma| = (p+q) - \#(\sigma)$.
  \item \textbf{Annular noncrossing set partitions $\NCA(p,q)$:} A partition $\pi \in \Pi(p+q)$ belongs to $\NCA(p,q)$ iff $\pi$ is the cycle partition of some permutation in $\SNC(p,q)$. Equivalently (cutting the annulus along a simple radial arc disjoint from the spanning trees of the blocks of $\pi$), $\pi \in \NCA(p,q)$ iff either (i) $\pi$ has no through-block and $\pi|_O \in \NC(p)$, $\pi|_I \in \NC(q)$, or (ii) $\pi$ has a through-block and $\pi$ is $\omega_{r_1,r_2}$-noncrossing for some cyclic rotations $r_1$ of $O$ and $r_2$ of $I$, where $\omega_{r_1,r_2} = (\mathrm{rot}_{r_1}(O), \mathrm{rot}_{r_2}(I))$.
  \item \textbf{Impermeable-annulus (non-through) family:} $\NCA^{\mathrm{nt}}(p,q) = \NC(p) \times \NC(q)$ ($p \ge q \ge 2$).
  \item \textbf{Rotational orbifold family $\NCB(p)$:} Partitions in $\NC(2p)$ invariant under the antipodal $180^\circ$ rotation $i \mapsto i + p \pmod{2p}$, representing noncrossing partitions on the $\mathbb{Z}_2$-quotient cone/annulus~\cite{BGN03}.
\end{enumerate}
\end{definition}

\begin{definition}[Canonical Boundary Families~\cite{SU91,Rei97,AR04}]\label{def:boundary}
The canonical planar noncrossing partition families equipped with a boundary structure are:
\begin{enumerate}
  \item \textbf{Arc-boundary family $\NC_\partial(m)$ ($m \ge 1$):} Partitions of $m$ vertices along a boundary arc of a disc together with one complementary boundary segment $\partial$, canonically isomorphic to $\NC(m+1)$.
  \item \textbf{Circular-boundary family $\NC_\partial^\circ(m)$ ($m \ge 1$):} Partitions $\pi \in \Pi(m+1)$ of $m$ vertices $\{0,\dots,m-1\}$ on an interior circle surrounded by an outer boundary circle $\partial = m$, where interior chords do not cross and any subset of exposed vertices may attach radially to $\partial$; combinatorially, $\pi \in \NC_\partial^\circ(m) \iff \pi|_{[m]} \in \NC(m)$.
  \item \textbf{Half-disc reflection-boundary families $\NCsym(m)$ ($m \ge 2$):} Sublattices of $\NC(m)$ fixed by a reflection involution across a diameter of the $m$-gon (passing between vertex pairs as $\NCsym_{\mathrm{edge}}(2k)$, through two opposite vertices as $\NCsym_{\mathrm{vert}}(2k)$, or through one vertex and an opposite edge for odd $m$)~\cite{SU91,Sim00}.
  \item \textbf{Zero-block (central boundary) Type-$B$ and Type-$D$ lattices:} Reiner's $\NCB(m)$ ($m \ge 1$, at most one symmetric zero/boundary block)~\cite{Rei97} and Athanasiadis--Reiner's $\NCD(m)$ ($m \ge 2$)~\cite{AR04}.
\end{enumerate}
\end{definition}

\section{Incidence Algebra Foundations and Annular Lattice Dichotomy}

\subsection{Incidence Algebra Symmetry and Isomorphism Invariance}

\begin{proposition}[Existence, Uniqueness, Dual Symmetry, and Invariance]\label{prop:mobius-basic}
Let $P$ be any finite poset.
\begin{enumerate}
  \item Equation~\eqref{eq:mobius-def} uniquely determines $\mu_P : P \times P \to \mathbb{Z}$, and in the incidence algebra of $P$ one has both $\mu_P \star \zeta_P = \delta$ and $\zeta_P \star \mu_P = \delta$.
  \item On the order dual $P^{\mathrm{op}}$, $\mu_{P^{\mathrm{op}}}(x, y) = \mu_P(y, x)$ for all $x, y \in P$; hence $\mathcal{S}_{\mathrm{full}}(P^{\mathrm{op}}) = \mathcal{S}_{\mathrm{full}}(P)$, while $\mathcal{S}_-(P^{\mathrm{op}}) = \mathcal{S}_+(P)$.
  \item Any poset (or lattice) isomorphism $\varphi : P \xrightarrow{\sim} Q$ satisfies $\mu_Q(\varphi(x), \varphi(y)) = \mu_P(x, y)$ for all $x, y \in P$, and therefore preserves $\mathcal{S}_-$, $\mathcal{S}_+$, and $\mathcal{S}_{\mathrm{full}}$.
\end{enumerate}
\end{proposition}
\begin{proof}
Constructing $\mu_P(x,y)$ by strong induction on the strict subinterval cardinality $|\{z \in P : x \le z < y\}|$ gives a left inverse $\mu_P \star \zeta_P = \delta$. Applying the same construction to $P^{\mathrm{op}}$ produces $r(x,y) = \mu_{P^{\mathrm{op}}}(y,x)$ satisfying $\zeta_P \star r = \delta$. By associativity of incidence-algebra convolution, any left inverse $\tilde{\mu}$ of $\zeta_P$ satisfies $\tilde{\mu} = \tilde{\mu} \star (\zeta_P \star r) = (\tilde{\mu} \star \zeta_P) \star r = r$. Hence $\mu_P = r$ is unique, establishing both (1) and (2). Pulling $\mu_Q$ back along $\varphi$ satisfies~\eqref{eq:mobius-def} on $P$, so uniqueness gives $\mu_Q(\varphi(x),\varphi(y)) = \mu_P(x,y)$ and (3).
\end{proof}

\subsection{Complete Lattice Classification on the Annulus}

\begin{theorem}[Non-Lattice Property of $\SNC(p,q)$]\label{thm:snc-nonlattice}
For $p, q \ge 1$, the annular noncrossing permutation poset $\SNC(p,q)$ is a lattice if and only if $p = q = 1$ (where $\SNC(1,1) \cong C_2$). For all $p + q \ge 3$, $\SNC(p,q)$ is neither a join-semilattice nor a meet-semilattice.
\end{theorem}
\begin{proof}
For $p = q = 1$, $\SNC(1,1) = \{\mathrm{id}, (0,1)\} \cong C_2$. Suppose $p + q \ge 3$; by symmetry between outer and inner circles we may assume $p \ge 2$ and $q \ge 1$. Consider the transpositions $\tau_1 = (0,1)$ (non-through) and $\tau_2 = (0, p)$ (through), together with the $3$-cycles $\gamma_1 = (0, 1, p)$ and $\gamma_2 = (0, p, 1)$, all fixing $[p+q] \setminus \{0, 1, p\}$. Direct cycle counting against $\gamma_{p,q}$ shows $\tau_1, \tau_2, \gamma_1, \gamma_2 \in \SNC(p,q)$. Since $|\tau_i| = 1$, $|\gamma_j| = 2$, and $\tau_i^{-1}\gamma_j$ is a transposition on $\{0, 1, p\}$ of reflection length $1$ for every $i, j \in \{1, 2\}$, we have $\tau_i \le \gamma_j$ for all $i, j \in \{1,2\}$. Because $|\gamma_1| = |\gamma_2| = 2$ and $\gamma_1 \ne \gamma_2$, any common upper bound $u$ of $\tau_1, \tau_2$ below both $\gamma_1$ and $\gamma_2$ would satisfy $|u| \le 1$, impossible since $\tau_1 \ne \tau_2$ both have length $1$. Dually, $\gamma_1$ and $\gamma_2$ have three maximal common lower bounds $(0,1)$, $(0,p)$, and $(1,p)$, failing both unique joins and unique meets.
\end{proof}

\begin{theorem}[Punctured-Annulus Equivalence and Enumeration]\label{thm:nca-p1}
For every $p \ge 1$:
\begin{enumerate}
  \item $\NCA(p, 1) = \{\pi \in \Pi(p+1) : \pi|_{[p]} \in \NC(p)\} = \NC_\partial^\circ(p)$, and $\NCA(p, 1)$ is a finite graded lattice (closed under $\wedge_{\Pi(p+1)}$ with $\hat{1}$).
  \item $|\NCA(p, 1)| = \frac{p+3}{2} C_p = \binom{2p}{p} - \binom{2p-1}{p+1}$, and $\mu_{\NCA(p,1)}(\hat{0}, \hat{1}) = (-1)^p \binom{2p-2}{p-1}$.
  \item $\NCA(p, 1)$ has $\binom{p+1}{2}$ atoms and $p^2 - p + 1$ coatoms; since $(p^2 - p + 1) - \binom{p+1}{2} = \binom{p-1}{2} > 0$ for all $p \ge 3$, $\NCA(p, 1)$ is non-self-dual for every $p \ge 3$.
\end{enumerate}
\end{theorem}
\begin{proof}
(1) If $\pi \in \NCA(p,1)$, either $\{p\} \in \pi$ and $\pi|_{[p]} \in \NC(p)$ by definition, or $\pi$ is $\omega_{r_1,0}$-noncrossing for some cyclic shift $\mathrm{rot}_{r_1}$ of $[p]$; deleting index $p$ preserves noncrossingness on $\mathrm{rot}_{r_1}([p])$, and cyclic invariance of $\NC(p)$ gives $\pi|_{[p]} \in \NC(p)$. Conversely, let $\pi_0 = \pi|_{[p]} \in \NC(p)$. If $\{p\} \in \pi$, then $\pi \in \NCA(p,1)$. Otherwise, since $p$ is a single element, $p$ belongs to $B \cup \{p\}$ for a unique block $B \in \pi_0$. Choose any $i_0 \in B$ and rotate $[p]$ by $r_1 = i_0 + 1$ steps so that $i_0$ sits at position $p-1$, immediately adjacent to position $p$ (which holds element $p$). Suppose $a < b < c < d$ crossed in $\omega_{r_1,0}$. Necessarily $d = p$, so $\omega(b), \omega(p) \in B \cup \{p\}$, which implies $\omega(p-1) = i_0 \in B$ lies in the same block as $\omega(b)$ and $\omega(p)$. Since $\omega(c)$ lies in the distinct block of $\omega(a)$, we have $c \ne p-1$, hence $a < b < c < p-1$, producing a crossing of $\pi_0$ on $\mathrm{rot}_{r_1}([p])$---contradiction. Because $(\sigma \wedge_\Pi \tau)|_{[p]} = \sigma|_{[p]} \wedge_\Pi \tau|_{[p]}$ and $\NC(p)$ is closed under $\wedge_\Pi$, $\NCA(p,1)$ is a meet-subsemilattice of $\Pi(p+1)$ containing $\hat{0}$ and $\hat{1}$, hence a finite lattice.

(2) Each $\pi_0 \in \NC(p)$ has $1 + |\pi_0|$ lifts in $\NCA(p,1)$ (either $\{p\}$ is a singleton or $p$ joins one of the $|\pi_0|$ blocks of $\pi_0$). By Kreweras self-duality of $\NC(p)$, the average block count over $\NC(p)$ is $\frac{p+1}{2}$, yielding $|\NCA(p,1)| = C_p(1 + \frac{p+1}{2}) = \frac{p+3}{2}C_p$. For $\mu_{\NCA(p,1)}(\hat{0},\hat{1})$, apply Weisner's theorem in $L = \NCA(p,1)$ to the atom $e = \{\{0,p\},\{1\},\dots,\{p-1\}\}$: $\mu_L(\hat{0},\hat{1}) = -\sum_{x \ne \hat{1},\, x \vee_L e = \hat{1}} \mu_L(\hat{0},x)$. For $\{p\} \in x$, $x \vee_L e = \hat{1}$ holds uniquely for $x = \{\{0,\dots,p-1\},\{p\}\}$, contributing $(-1)^{p-1}C_{p-1}$. For $p$ merged into a block $B_2 \ne B_1 \ni 0$ of $\pi_0 = x|_{[p]} \in \NC(p)$, $x \vee_L e = \hat{1}$ iff merging $B_1$ and $B_2$ in $\NC(p)$ yields $\hat{1}_{\NC(p)}$; summing the factorized lower intervals $[\hat{0},x] \cong \NCA(|B_2|,1) \times \prod_{C \ne B_2} \NC(|C|)$ yields the generating-function relation $A(z) = \frac{z}{\sqrt{1-4z}} = \sum_{p \ge 1} \binom{2p-2}{p-1} z^p$ for $a_p = (-1)^p \mu_{\NCA(p,1)}(\hat{0},\hat{1})$.

(3) Every $2$-element subset of $[p+1]$ restricts on $[p]$ to a partition with block sizes $\le 2$ (automatically noncrossing), giving $\binom{p+1}{2}$ atoms. A coatom has two blocks $\{A^c, A \cup \{p\}\}$ with $\{A^c, A\} \in \NC(p)$ (or $A = \emptyset$); the nonempty proper subsets $A^c \subsetneq [p]$ such that $\{A^c, A\} \in \NC(p)$ are precisely the $p(p-1)$ proper cyclic intervals of $[p]$, giving $1 + p(p-1) = p^2 - p + 1$ coatoms.
\end{proof}

\begin{theorem}[Complete Lattice Dichotomy for $\NCA(p,q)$]\label{thm:nca-dichotomy}
Let $p \ge q \ge 1$. The annular set-partition poset $\NCA(p,q)$ under refinement is a lattice if and only if $q = 1$ or $(p,q) \in \{(2,2), (3,2)\}$. Specifically:
\begin{enumerate}
  \item For $(p,q) \in \{(1,1), (2,1), (2,2), (3,1), (3,2)\}$, $\NCA(p,q) = \Pi(p+q)$.
  \item For $q = 1$ and $p \ge 4$, $\NCA(p,1) \subsetneq \Pi(p+1)$ is a proper lattice by Theorem~\ref{thm:nca-p1}.
  \item For all $p \ge q \ge 2$ with $p + q \ge 6$, $\NCA(p,q)$ fails both unique joins and unique meets.
\end{enumerate}
\end{theorem}
\begin{proof}
Items (1) and (2) are immediate for $p+q \le 4$; for $(3,2)$, any partition of $5$ elements with no through-block restricts to subsets of sizes $\le 3$ (always noncrossing), and any partition with a through-block has $\le 2$ blocks of size $\ge 2$, one of which can be placed on an adjacent outer--inner cut pair $[2,3]$ by rotating its outer and inner endpoints to indices $2$ and $3$, preventing any crossing.

For (3), either $p \ge 4, q \ge 2$ or $p \ge 3, q \ge 3$:
\begin{itemize}
  \item \emph{Case $p \ge 4, q \ge 2$:} Define $\sigma = \{\{3, p\}\}$ and $\tau = \{\{0, 2\}, \{1, p+1\}\}$ (with all unmentioned elements of $[p+q]$ as singletons), and let
  \[
    u_1 = \bigl\{\{0, 1, 2, p+1\},\, \{3, p\}\bigr\}, \qquad u_2 = \bigl\{\{0, 2, 3, p\},\, \{1, p+1\}\bigr\}.
  \]
  On the unrotated cut $\omega_{0,0} = \mathrm{id}$, the block $\{3, p\}$ of $u_1$ lies in the index interval $[3, p]$ while $\{0, 1, 2, p+1\}$ lies strictly outside $[3, p]$, so $\sigma, u_1 \in \NCA(p,q)$. On the rotated cut $\omega_{2,1}$, vertex $1$ sits at index $p-1$ and vertex $p+1$ sits at index $p$, placing $\{1, p+1\}$ in the adjacent pair $[p-1, p]$ disjoint from all other blocks, so $\tau, u_2 \in \NCA(p,q)$. Now suppose $u \in \NCA(p,q)$ were a least upper bound of $\sigma$ and $\tau$. Then $\sigma, \tau \le u \le u_1, u_2$ forces $u = u_1 \wedge_\Pi u_2 = \{\{0,2\}, \{1,p+1\}, \{3,p\}\} =: w_0$. Since $w_0$ has through-blocks, $w_0 \in \NCA(p,q)$ would give a cut $\omega_{r_1,r_2}$ on which $w_0$ is noncrossing. In $\mathrm{rot}_{r_1}(O)$, the four outer vertices $(0,1,2,3)$ appear at indices $k_0, k_1, k_2, k_3 < p$ in one of the four cyclic orders, while $p, p+1$ appear at indices $m_0, m_1 \ge p$. If $k_0 < k_1 < k_2$, then $k_0 < k_1 < k_2 < m_1$ crosses $\{0,2\}$ with $\{1,p+1\}$; in the remaining two cyclic shifts $(k_1 < k_2 < k_3 < k_0)$ and $(k_2 < k_3 < k_0 < k_1)$, we have $k_2 < k_3 < k_0 < m_0$, crossing $\{2,0\}$ with $\{3,p\}$. Thus $w_0 \notin \NCA(p,q)$, a contradiction.
  \item \emph{Case $p \ge 3, q \ge 3$:} Define $\sigma = \{\{2, p\}\}$, $\tau = \{\{0, p+1\}, \{1, p+2\}\}$, $u_1 = \{\{0, 1, p+1, p+2\}, \{2, p\}\}$, and $u_2 = \{\{0, 2, p, p+1\}, \{1, p+2\}\}$. By the same interval-separation test on $\omega_{0,0}$ (interval $[2,p]$) and $\omega_{2,2}$ (adjacent pair $[p-1, p]$), we have $\sigma, \tau, u_1, u_2 \in \NCA(p,q)$, and any least upper bound would equal $w_0 = u_1 \wedge_\Pi u_2 = \{\{0,p+1\}, \{1,p+2\}, \{2,p\}\}$. In any cut $\omega_{r_1,r_2}$, $(0,1,2)$ appear at $k_0, k_1, k_2 < p$ in one of three cyclic orders and $(p, p+1, p+2)$ appear at $m_0, m_1, m_2 \ge p$ in one of three cyclic orders; in every one of the $3 \times 3 = 9$ cyclic cases, two of the three bridges $\{0,p+1\}, \{1,p+2\}, \{2,p\}$ satisfy $k_a < k_b < m_{\phi(a)} < m_{\phi(b)}$ and cross. Hence $w_0 \notin \NCA(p,q)$.
\end{itemize}
Since every $\NCA(p,q)$ contains $\hat{0}$ and $\hat{1}$, any finite meet-semilattice with $\hat{1}$ would automatically have unique joins ($x \vee y = \bigwedge \{w : x, y \le w\}$); therefore failure of unique joins implies failure of unique meets as well.
\end{proof}

\section{Universal Canonical Rigidity and Disproof}

\subsection{Resolution of the \texorpdfstring{$\binom{2n}{n}$}{C(2n, n)} Equinumerosity Family}

\begin{theorem}[Rotational vs.\ Reflection-Boundary Separation Failure]\label{thm:equinumerous}
For every $n \ge 1$, $|\NCB(n)| = |\NCsym_{\mathrm{edge}}(2n)| = |\NCsym_{\mathrm{vert}}(2n)| = \binom{2n}{n}$.
\begin{enumerate}
  \item At $n = 1$, all three lattices are isomorphic to $C_2$.
  \item For all $n \ge 2$, the three lattices are pairwise non-isomorphic; Kreweras complementation induces an order-duality isomorphism $(\NCsym_{\mathrm{edge}}(2n))^{\mathrm{op}} \cong \NCsym_{\mathrm{vert}}(2n)$; and
  \[
    \bigl|\mu_{\NCB(n)}(\hat{0},\hat{1})\bigr| = \binom{2n-1}{n-1} > C_{n-1} = \bigl|\mu_{\NCsym_{\mathrm{edge}}(2n)}(\hat{0},\hat{1})\bigr| = \bigl|\mu_{\NCsym_{\mathrm{vert}}(2n)}(\hat{0},\hat{1})\bigr|.
  \]
  Hence $\NCB(n)$ disagrees with both reflection-boundary lattices on both $\mathcal{S}_-$ and $\mathcal{S}_{\mathrm{full}}$ (for example at $n = 2$: $\mathcal{S}_-(\NCB(2)) = \{-1^4, 1^1, 3^1\}$ vs.\ $\{-1^3, 1^3\}$ and $\{-1^2, 0^2, 1^2\}$). Moreover, $\mathcal{S}_-(\NCsym_{\mathrm{edge}}(2n)) \ne \mathcal{S}_-(\NCsym_{\mathrm{vert}}(2n))$ even between the two boundary axes.
\end{enumerate}
\end{theorem}
\begin{proof}
Because $180^\circ$ rotation and diameter reflections are lattice automorphisms of $\NC(2n)$, their fixed-point subposets are sublattices of $\NC(2n)$. By Reiner~\cite{Rei97} and Simion~\cite{Sim00}, all three have cardinality $\binom{2n}{n}$, with $\mu_{\NCB(n)}(\hat{0},\hat{1}) = (-1)^n\binom{2n-1}{n-1}$ and $\mu_{\NCsym(2n)}(\hat{0},\hat{1}) = (-1)^n C_{n-1}$. For $n \ge 2$, we have $n C_{n-1} = \binom{2n-2}{n-1} < \binom{2n-1}{n-1}$, whence $\binom{2n-1}{n-1} > n C_{n-1} \ge 2 C_{n-1} > C_{n-1}$. Because every proper interval in $\NCB(n)$ factors into smaller $\NCB$ and $\NC$ lattices, $\binom{2n-1}{n-1}$ is also the unique maximum of $|\mu(x,y)|$ in $\mathcal{S}_{\mathrm{full}}(\NCB(n))$, whereas $\max|\mu| = C_{n-1}$ in $\mathcal{S}_{\mathrm{full}}(\NCsym(2n))$. Finally, Kreweras complementation rotates vertex adjacency by half a step, conjugating an edge-axis reflection into a vertex-axis reflection while reversing order; at $n \ge 2$, $\NCsym_{\mathrm{edge}}(2n)$ has no zero-block-type pinch above atoms whereas $\NCsym_{\mathrm{vert}}(2n)$ has coatoms with $\mu(\hat{0}, x) = 0$ (e.g.\ $0^0$ vs.\ $0^2$ at $n=2$).
\end{proof}

\subsection{Universal All-Parameter Rigidity Theorem}

\begin{theorem}[Universal Canonical Rigidity]\label{thm:universal-rigidity}
Let $\mathcal{A}$ range over all canonical annular noncrossing partition lattices---$\NCA(p,1)$ ($p \ge 1$), $\NCA(2,2) = \Pi_4$, $\NCA(3,2) = \Pi_5$, $\NC(p)\times\NC(q)$ ($p \ge q \ge 2$), and rotational $\NCB(p)$ ($p \ge 1$)---and let $\mathcal{B}$ range over all canonical boundary noncrossing partition lattices---$\NC_\partial(m) \cong \NC(m+1)$ ($m \ge 1$), $\NC_\partial^\circ(m)$ ($m \ge 1$), $\NCB(m)$ ($m \ge 1$), $\NCD(m)$ ($m \ge 2$), and $\NCsym(m)$ ($m \ge 2$). Then $(|\mathcal{A}|, |\mu_{\mathcal{A}}(\hat{0},\hat{1})|) = (|\mathcal{B}|, |\mu_{\mathcal{B}}(\hat{0},\hat{1})|)$ holds if and only if $(\mathcal{A}, \mathcal{B})$ is one of the six canonical isomorphic pairs:
\begin{enumerate}
  \item $|\mathcal{A}| = |\mathcal{B}| = 2$ ($C_2$): $\mathcal{A} \in \{\NCA(1,1), \NCB(1)\}$ and $\mathcal{B} \in \{\NC_\partial(1),\allowbreak\, \NC_\partial^\circ(1),\allowbreak\, \NCB(1),\allowbreak\, \NCsym(2)\}$;
  \item $|\mathcal{A}| = |\mathcal{B}| = 4$ ($B_2$): $\mathcal{A} = \NC(2)\times\NC(2)$ and $\mathcal{B} = \NCD(2)$;
  \item $|\mathcal{A}| = |\mathcal{B}| = 5$ ($\Pi_3$): $\mathcal{A} = \NCA(2,1)$ and $\mathcal{B} \in \{\NC_\partial(2), \NC_\partial^\circ(2)\}$;
  \item $|\mathcal{A}| = |\mathcal{B}| = 15$ ($\Pi_4$): $\mathcal{A} = \NCA(2,2)$ and $\mathcal{B} = \NC_\partial^\circ(3) = \NCA(3,1)$;
  \item $\mathcal{A} = \NCA(p,1)$ and $\mathcal{B} = \NC_\partial^\circ(p)$ for $p \ge 1$ (equal by Theorem~\ref{thm:nca-p1});
  \item $\mathcal{A} = \NCB(p)$ and $\mathcal{B} = \NCB(p)$ for $p \ge 1$.
\end{enumerate}
Consequently, \textbf{no} pair of a canonical annular family and a canonical boundary family forms non-isomorphic lattices with the same M\"obius function spectrum (under $\mathcal{S}_-$ or $\mathcal{S}_{\mathrm{full}}$).
\end{theorem}
\begin{proof}
Both $\mathcal{S}_-(L)$ and $\mathcal{S}_{\mathrm{full}}(L)$ determine $|L|$ and $|\mu_L(\hat{0},\hat{1})|$ across these families ($\sum_{v} \mathrm{mult}_{\mathcal{S}_-}(v) = |L|$ and $\max_{x \le y}|\mu_L(x,y)| = |\mu_L(\hat{0},\hat{1})|$ for $|L| > 2$, while $\mu_{\NCsym(2k+1)}(\hat{0},\hat{1}) = 0$ for odd $2k+1 \ge 3$, which never matches any canonical annular lattice since $|\mu_{\mathcal{A}}(\hat{0},\hat{1})| \ge 1$ everywhere). Writing $R(L) = |L| / |\mu_L(\hat{0},\hat{1})|$ whenever $\mu_L(\hat{0},\hat{1}) \ne 0$, the exact binomial recurrences yield the division-free identities:
\[
\begin{aligned}
  p(p+1)\,|\NCA(p,1)| &= (p+3)(2p-1)\,|\mu|, & R &= \tfrac{(p+3)(2p-1)}{p(p+1)} \in [2, 3) \; (R = 2 \Leftrightarrow p = 1), \\
  (m+2)\,|\NC_\partial(m)| &= 2(2m+1)\,|\mu|, & R &= \tfrac{4m+2}{m+2} < 4 \; (R \ge 3 \text{ for } m \ge 4), \\
  |\NCB(m)| &= 2\,|\mu|, & R &= 2, \\
  (3m-4)\,|\NCD(m)| &= 2(3m-2)\,|\mu|, & R &= \tfrac{2(3m-2)}{3m-4} \in (2, 4] \; (R < 4 \text{ for } m \ge 3), \\
  |\NCsym(2k+2)| &= 2(2k+1)\,|\mu|, & R &= 2(2k+1) \in \{2, 6, 10, 14, 18, \dots\}, \\
  (p+1)(q+1)\,|\NC(p){\times}\NC(q)| &= 4(2p-1)(2q-1)\,|\mu|, & R &\in [4, 16) \; (p \ge q \ge 2).
\end{aligned}
\]
Comparing $R(\mathcal{A}) = R(\mathcal{B})$ and monotonic cardinalities eliminates every non-isomorphic case:
\begin{itemize}
  \item \emph{$\mathcal{A} = \NCA(p,1)$:} Since $R(\mathcal{A}) < 3$, matching $\NC_\partial(m)$ forces $m \in \{1,2,3\}$ ($|\mathcal{B}| \in \{2,5,14\}$), yielding only $(p,m) \in \{(1,1),(2,2)\}$ ($|\NCA(3,1)| = 15 \ne 14$); matching $\NCB(m)$ ($R=2$) forces $(p+3)(2p-1)=2p(p+1) \implies p=1=m$; matching $\NCsym(2k+2)$ ($R = 2(2k+1) < 3$) forces $k=0$, $m=2, p=1$. Against $\NCD(m)$: for $m \in \{2,3,4\}$, $|\NCD(m)| \in \{4, 14, 50\} \notin \{2, 5, 15, 49, 168, \dots\}$; for $m \ge 5$, the identity $2m\,|\mu_{\NCD(m)}| = (3m-4)\binom{2m-2}{m-1}$ gives the strict central-binomial sandwich
  \[
    \binom{2m-2}{m-1} < |\mu_{\NCD(m)}(\hat{0},\hat{1})| < \binom{2m}{m},
  \]
  because $2m < 3m-4 < 4(2m-1)$ for all $m \ge 5$, so $|\mu_{\NCD(m)}|$ cannot equal $|\mu_{\NCA(p,1)}| = \binom{2p-2}{p-1}$ for any $p \ge 1$.
  \item \emph{$\mathcal{A} \in \{\NCA(2,2), \NCA(3,2)\}$:} $(|\mathcal{A}|, |\mu|) \in \{(15, 6), (52, 24)\}$ matches only $\NC_\partial^\circ(3)$ at $(15, 6)$.
  \item \emph{$\mathcal{A} = \NC(p)\times\NC(q)$ ($p \ge q \ge 2$):} Here $4 \le R(\mathcal{A}) < 16$, which immediately excludes $\NC_\partial(m)$ ($R < 4$), $\NC_\partial^\circ(m)$ ($R < 3$), $\NCB(m)$ ($R = 2$), and $\NCD(m)$ for $m \ge 3$ ($R \le 14/5 < 4$), leaving $(p,q,m) = (2,2,2)$ ($B_2 \cong \NCD(2)$). Against $\NCsym(2k+2)$, the bound $4 \le 2(2k+1) < 16$ restricts to $k \in \{1,2,3\}$, where $(|\mathcal{B}|, |\mu_{\mathcal{B}}|) \in \{(6,1), (20,2), (70,5)\}$. Meanwhile $|\mu_{\mathcal{A}}| = C_{p-1}C_{q-1} \le 5$ with $p \ge q \ge 2$ holds only for $(p,q) \in \{(2,2), (3,2), (4,2)\}$, whose cardinalities $C_p C_q \in \{4, 10, 28\}$ are disjoint from $\{6, 20, 70\}$.
  \item \emph{$\mathcal{A} = \NCB(p)$ ($R = 2$):} Matching $R = 2$ gives $m = 1, p = 1$ against $\NC_\partial$ and $\NC_\partial^\circ$, $p = m$ against $\NCB$, no solutions against $\NCD$ ($R > 2$), and $k = 0$ ($m = 2, p = 1$) against $\NCsym(2k+2)$. \qedhere
\end{itemize}
\end{proof}

\section{Scope Boundary: Arbitrary Induced Noncrossing Sublattices}

To confirm that Theorem~\ref{thm:universal-rigidity} reflects genuine rigidity of canonical surface families rather than a general poset-theoretic law, we examine what happens when one drops the canonical annular/boundary requirement and permits arbitrary finite subposet lattices inside $\NC(N)$.

\begin{lemma}[Pinch-Point Vanishing]\label{lem:pinch}
Let $L$ be a finite lattice containing a pinch element $m \in L$ comparable to every element of $L$ ($\forall z \in L,\, z \le m \text{ or } m \le z$). Then $\mu_L(x, y) = 0$ for all $x < m < y$.
\end{lemma}
\begin{proof}
Fix $x < m$ and induct on $|\{w : m \le w \le y\}|$ for $y > m$. In $\sum_{x \le z \le y} \mu_L(x,z) = 0$, every $z \in [x, y]$ satisfies either $x \le z \le m$ or $m < z \le y$. The sum over $x \le z \le m$ vanishes by~\eqref{eq:mobius-def} since $x < m$, and $\mu_L(x,z) = 0$ for $m < z < y$ by the induction hypothesis, forcing $\mu_L(x,y) = 0$.
\end{proof}

\begin{proposition}[Induced Sublattice Separation Inside $\NC(4)$ and $\NC(8)$]\label{prop:scope}
\leavevmode
\begin{enumerate}
  \item \textbf{One-variable separation in $\NC(4)$ (size $5$):} Let $L_{\mathrm{stem}}$ and $L_{N_5}$ be the $5$-element subposets of $\NC(4)$ on $\{0,1,2,3\}$ defined by
  \[
    L_{\mathrm{stem}} = \bigl\{\hat{0},\, \{2,3\},\, \{1,2\},\, \{1,2,3\},\, \hat{1}\bigr\}, \qquad
    L_{N_5} = \bigl\{\hat{0},\, \{2,3\},\, \{1,2\},\, \{\{0,1\},\{2,3\}\},\, \hat{1}\bigr\}.
  \]
  Both are lattices under refinement, $L_{\mathrm{stem}} \not\cong L_{N_5}$, and $\mathcal{S}_-(L_{\mathrm{stem}}) = \mathcal{S}_-(L_{N_5}) = \{1^2, -1^2, 0^1\}$.
  \item \textbf{Simultaneous all-spectrum separation in $\NC(8)$ (size $10$):} Embed $2^{\{0,\dots,6\}} \hookrightarrow \NC(8)$ by mapping each gap bitmask $S \subseteq \{0,\dots,6\}$ to the contiguous interval partition of $\{0,\dots,7\}$ merging $(g, g+1)$ iff $g \in S$. Let $L_{10}^{(A)}$ and $L_{10}^{(B)}$ correspond to the bitmask families
  \[
    \mathcal{M}_A = \{0, 1, 3, 7, 11, 15, 31, 47, 63, 127\}, \qquad
    \mathcal{M}_B = \{0, 1, 3, 5, 7, 15, 31, 47, 63, 127\},
  \]
  realizing $C_2 \triangledown A \triangledown A^{\mathrm{op}} \triangledown C_2$ and $C_2 \triangledown A^{\mathrm{op}} \triangledown A \triangledown C_2$ where $A$ is the $5$-element lattice $\hat{0} < u < \{v_1,v_2\} < \hat{1}$. Both $L_{10}^{(A)}$ and $L_{10}^{(B)}$ are self-dual graded sublattices of $\NC(8)$ of rank $7$ with identical $\mathcal{S}_- = \mathcal{S}_+ = \{1^1, -1^1, 0^8\}$, identical $\chi(t) = t^7 - t^6$, and identical $\mathcal{S}_{\mathrm{full}} = \{1^{12}, -1^{11}, 0^{30}\}$, yet $L_{10}^{(A)} \not\cong L_{10}^{(B)}$ (rank $2$ has $1$ element in $L_{10}^{(A)}$ vs.\ $2$ elements in $L_{10}^{(B)}$).
\end{enumerate}
\end{proposition}
\begin{proof}
Part (1) follows by direct computation on the two $5$-element Hasse diagrams ($L_{\mathrm{stem}}$ has $\mathcal{S}_+ = \{1^1, -1^1, 0^3\}$ whereas $L_{N_5}$ is self-dual with $\mathcal{S}_+ = \{1^2, -1^2, 0^1\}$). For Part (2), both $\mathcal{M}_A$ and $\mathcal{M}_B$ are closed under bitwise $\cap$ and $\cup$ (hence sublattices of $2^{\{0,\dots,6\}} \subseteq \NC(8)$). By Lemma~\ref{lem:pinch} applied at the pinch masks $1$, $15$, and $63$, every interval crossing a pinch has $\mu = 0$; the remaining proper intervals in both lattices consist of the same multiset $\{C_2, A, A^{\mathrm{op}}, C_2\}$ up to order duality, so Proposition~\ref{prop:mobius-basic}(2) yields $\mathcal{S}_{\mathrm{full}}(L_{10}^{(A)}) = \mathcal{S}_{\mathrm{full}}(L_{10}^{(B)})$ alongside $\mathcal{S}_- = \mathcal{S}_+ = \{1^1, -1^1, 0^8\}$. Non-isomorphism holds since $|\{x \in L_{10}^{(A)} : |(\downarrow x)| = 3\}| = 1 \ne 2 = |\{x \in L_{10}^{(B)} : |(\downarrow x)| = 3\}|$.
\end{proof}

\section{Formal Verification in Core Lean 4.19.0 and Python}

All definitions and theorems are formalized from first principles in Core Lean~4 (toolchain \texttt{leanprover/\allowbreak lean4:\allowbreak v4.19.0}, compiled with \texttt{-DwarningAsError=true} and zero external dependencies) under \texttt{lean/\allowbreak Conjecture00000004770/}, with zero \texttt{sorry}, \texttt{admit}, \texttt{native\_decide}, or custom axioms:
{\sloppy\hbadness=2000
\begin{itemize}
  \item \texttt{PosetMobius.lean}: Finite posets, lattice axioms, order and lattice isomorphisms, incidence-algebra convolution, existence and uniqueness of $\mu_P$ (\texttt{mobius\_spec}, \texttt{mobius\_unique}), M\"obius inversion (\texttt{mobius\_\allowbreak inversion}), dual symmetry (\texttt{mobius\_dual}), and isomorphism invariance of $\mathcal{S}_-$, $\mathcal{S}_+$, and $\mathcal{S}_{\mathrm{full}}$ (Proposition~\ref{prop:mobius-basic}).
  \item \texttt{PartitionsAndAnnulus.lean}: Set partitions, disc $\NC(N)$, cut-annulus $\NCA(p,q)$, universal punctured-annulus equivalence $\NCA(p,1) = \NC_\partial^\circ(p)$ (\texttt{annularNC\_\allowbreak p\_\allowbreak one\_\allowbreak iff\_\allowbreak circBoundary}, Theorem~\ref{thm:nca-p1}), and universal join-semilattice failure for all $p \ge q \ge 2$ with $p+q \ge 6$ (\texttt{annularNC\_\allowbreak no\_\allowbreak lub\_\allowbreak of\_\allowbreak pq\_\allowbreak ge\_\allowbreak two\_\allowbreak sum\_\allowbreak ge\_\allowbreak six}, Theorem~\ref{thm:nca-dichotomy}).
  \item \texttt{CanonicalRigidity.lean}: Concrete size-$6$ disproofs (\texttt{snc21\_\allowbreak not\_\allowbreak lattice}, \texttt{size6\_\allowbreak trap\_\allowbreak resolution}, Theorems~\ref{thm:snc-nonlattice} and~\ref{thm:equinumerous}), division-free binomial and Catalan recurrences, the Type-$D$ central-binomial sandwich (\texttt{ncDMobius\allowbreak Abs\_\allowbreak strictly\_\allowbreak between\_\allowbreak central\allowbreak Binom}), and the all-parameter classification (\texttt{canonical\_\allowbreak invariants\_\allowbreak match\_\allowbreak iff\_\allowbreak isomorphic}, Theorem~\ref{thm:universal-rigidity}).
  \item \texttt{SubfamilyScope.lean} and root \texttt{Conjecture00000004770.lean}: Induced sublattice pairs in $\NC(4)$ and $\NC(8)$ (Proposition~\ref{prop:scope}) and top-level exports \texttt{conjecture\_\allowbreak 00000004770\_\allowbreak positive\_\allowbreak foundations} and \texttt{conjecture\_\allowbreak 00000004770\_\allowbreak disproof}.
\end{itemize}
}
The standalone script \texttt{verify.py} independently cross-checks $\SNC(p,q)$ and $\NCA(p,q)$ ($p+q \le 6$), the $\binom{2n}{n}$ trap ($n \le 4$), $\NCD(n)$ ($n \le 4$), $329{,}096$ canonical parameter pairs up to $50$, and both scope-boundary sublattice constructions.

\begin{thebibliography}{9}
\bibitem{AR04}
C.~A. Athanasiadis and V.~Reiner, \emph{Noncrossing partitions for the group $D_n$}, SIAM J. Discrete Math. \textbf{18} (2004), no.~2, 397--417.

\bibitem{BGN03}
P.~Biane, F.~Goodman, and A.~Nica, \emph{Non-crossing cumulants of type B}, Trans. Amer. Math. Soc. \textbf{355} (2003), no.~6, 2263--2303.

\bibitem{Kre72}
G.~Kreweras, \emph{Sur les partitions non crois\'ees d'un cycle}, Discrete Math. \textbf{1} (1972), no.~4, 333--350.

\bibitem{MN04}
J.~A. Mingo and A.~Nica, \emph{Annular non-crossing permutations and partitions, and second-order asymptotics for random matrices}, J. Reine Angew. Math. \textbf{571} (2004), 13--68.

\bibitem{NS06}
A.~Nica and R.~Speicher, \emph{Lectures on the Combinatorics of Free Probability}, London Math. Soc. Lecture Note Ser., vol.~335, Cambridge Univ. Press, 2006.

\bibitem{Rei97}
V.~Reiner, \emph{Non-crossing partitions for classical reflection groups}, Discrete Math. \textbf{177} (1997), no.~1--3, 195--222.

\bibitem{Sim00}
R.~Simion, \emph{Noncrossing partitions}, Discrete Math. \textbf{217} (2000), no.~1--3, 367--409.

\bibitem{SU91}
R.~Simion and D.~Ullman, \emph{On the structure of the lattice of noncrossing partitions}, Discrete Math. \textbf{98} (1991), no.~3, 193--206.
\end{thebibliography}

\end{document}
